\subsection{A lattice obstruction to ordinary interlacing alone}
\begin{finding}[\fid{F18}: homogeneous difference constraints and a rectangular parity ray]\label{f:saturation}
Let $u$ range over an integral label set closed under positive-integer scaling. Fix a number of coordinates and a finite list of inequalities
\[
 P(u)=\{x\in\mathbb R^s:x_i-x_j\le b_{ij}(u)\},\qquad x_0=0,
\]
where every specified $b_{ij}(u)$ is integral and $b_{ij}(Nu)=Nb_{ij}(u)$ for every $N\ge1$. Assume $P(u)$ is bounded, and put $h(u)=\#(P(u)\cap\Z^s)$. Then
\[
 h(Nu)>0\quad\Longleftrightarrow\quad h(u)>0.
\]
The same property holds for a fixed finite sum of such counting functions. Thus these models cannot count all two-row rectangular Kronecker coefficients. On the particular ray with all three partitions $(d,d)$, the established two-row formula gives
\[
 g((d,d),(d,d),(d,d))=\boldsymbol1_{2\mid d}
 =\#\bigl(d\{1/2\}\cap\Z\bigr)\qquad(d\ge1).
\]
This ray has an exact rational Ehrhart model despite its failure of the displayed ordinary support property.
\end{finding}
\begin{proof}
Homogeneity gives $P(Nu)=NP(u)$. If $x_i-x_j\le b$ with $b\in\Z$, then
$\lfloor x_i\rfloor\le\lfloor x_j+b\rfloor=\lfloor x_j\rfloor+b$.
Hence coordinatewise flooring, including $x_0=0$, takes every real point of $P(u)$ to an integral one. If $z\in P(Nu)\cap\Z^s$, apply this to $z/N$. Conversely an integral $x\in P(u)$ gives $Nx\in P(Nu)\cap\Z^s$. Positivity of a fixed finite sum is positivity of at least one summand, so the same argument applies. Boundedness is only needed to make these finite counts.

For $u=((1,1),(1,1),(1,1))$, the triple tensor product is the sign representation of $S_2$, so $g(u)=0$. The explicit $S_4$ character calculation following \fid{F13} gives $g(2u)=1$: the class sizes are $(1,6,3,8,6)$ and the character $(2,0,2,-1,0)$ has cubic average $(8+3\cdot8-8)/24=1$. Thus an exact model of the stated form on all two-row rectangles is impossible.

For the whole ray use IV, \S18.3, equation~(172), printed p.~128. In its column convention, $(d,d)=[d,0]$ as a two-row label and $[0,0,d,0]$ as a four-row label. The coefficient for the second rectangle is the constant term of that equation's generating function. Only $k=0$ and $j=0$ in its first sum survive; the second sum is empty. The remaining condition is that $n_3,n_2,n_1$ are all even and nonnegative. Substituting $(n_3,n_2,n_1)=(0,d,0)$ gives exactly the displayed parity formula. The single rational point $\{1/2\}$ has an integral point after dilation by $d$ exactly when $d$ is even.
\end{proof}
The difference-constraint result includes ordinary interlacing with integral endpoint rows. Its exclusions are essential: extra congruences, sum equations, nonhomogeneous endpoint shifts, a changing constraint list or a changed lattice can evade it. For example $2z=d$ retains precisely the parity that ordinary integral differences lose. The example concerns all two-row rectangles, not a model restricted only to IP's large-parameter region.

Nor does it contradict VI's definition of ``strictly saturated'' (Definition~1.2.4, printed p.~14). That definition tests each nonzero residue polynomial separately; here $f_{\rm even}=1$ and $f_{\rm odd}=0$, so its saturation index is zero. We use the displayed support property to avoid conflating the two notions. IV, Theorem~18.2 and equations~(163)--(164), explicitly retain congruences and remove columns in pairs. A successful theta comparison must preserve these arithmetic data or equivalent information; the source does not identify them with a particular GGP local sign.

In Grothendieck's language, the tensor group-scheme map $\GL_{2,\Z}\times\GL_{2,\Z}\to\GL_{4,\Z}$ induces, over $\Q$, a positive restriction map on representation $K_0$. The coefficient of $[\det^d\boxtimes\det^d]$ in the restriction of $[S_{(d,d)}(V\otimes W)]$ is this parity function. Thus an effective $K_0$ homomorphism does not by itself give the ordinary support property. The proof above is already the exact integral-lattice dictionary; its group-scheme framework is SGA~3 I, 4.7.2--3, with cote~90, scan~17, as an action/comodule anchor.
